\begin{frame}{기하학적 성격: 보로노이 다이어그램}
  \begin{itemize}
    \item \kb{$k=1$ kNN의 결정경계 = \kb{보로노이 다이어그램}의 경계}
      (Voroni, 1908)과 정확히 일치
    \item \kb{보로노이 셀}: 특정 훈련점 $x_i$에 가장 가까운 모든
      지점들의 영역 $\{x : d(x, x_i) \le d(x, x_j)\ \forall j\}$
      --- 각 $x_i$를 중심으로 한 \textbf{볼록 다각형}, 이웃 셀 경계는
      두 점의 \textbf{수직등분선}
  \end{itemize}
  \vskip0.8em
  \begin{block}{두 가지 실용적 의미}
    \begin{enumerate}\small
      \item $k=1$ 경계는 \textbf{항시 각진(비선형) 다각형의 합} ---
        아무리 많은 점을 써도 \textbf{곡선을 그을 수 없다}(충분히
        작은 직선 조각으로만 근사)
      \item \textbf{두 훈련점이 아주 가까우면} 그 수직등분선이
        두 점 사이를 지나, 근처 경계가 두 점의 \textbf{상대적 위치
       에만} 달려 있다 = 분산 얘기의 기하학적 번역
    \end{enumerate}
  \end{block}
  $k>1$에서는 \kb{$k$-보로노이}가 경계를 정의 --- ``$k$-이웃의
  클래스 분포가 50:50으로 바뀌는 곳''. $k$가 커질수록 매끄러워진다는
  직감의 기하학적 근거.
\end{frame}
